\section{Context Growth, Reconstruction Cost, and State Economics}
\label{sec:cost}

RaS has an economic motivation, but its central intellectual claim is continuity placement rather than a specific price forecast. A provider-neutral model is therefore preferable to assumptions about internal GPU scheduling, proprietary caches, margins, or subscription economics.

\subsection{A simple cumulative-context model}

Let $B$ be fixed base context, $g$ the average number of new historical tokens introduced per sequential interaction, and $n$ the number of reasoning steps. Under a naive full-history replay model, the input presented at step $i$ is
\begin{equation}
T_i = B + g(i-1).
\end{equation}
The cumulative supplied context is
\begin{align}
T_{\mathrm{persistent}}(n)
&=
\sum_{i=1}^{n}[B+g(i-1)] \\
&=
nB + \frac{gn(n-1)}{2}.
\label{eq:persistent-context}
\end{align}
For fixed $g>0$,
\begin{equation}
T_{\mathrm{persistent}}(n)\in\Theta(n^2).
\end{equation}

Now let $K_i$ be the repository/task context reconstructed for RaS step $i$. If
\begin{equation}
K_i \le K
\end{equation}
for bounded $K$, then
\begin{equation}
T_{\mathrm{RaS}}(n)
=
\sum_{i=1}^{n}K_i
\le nK,
\end{equation}
and, given a non-degenerate lower bound, cumulative supplied context is $\Theta(n)$.

Figure~\ref{fig:context-growth} illustrates only these mathematical assumptions.

\begin{figure}[htbp]
\centering
\input{figures/context-growth}
\caption{Conceptual cumulative-context growth under the stated models. The curves are illustrative mathematical shapes, not measured agent or provider costs.}
\label{fig:context-growth}
\end{figure}

This result must not be confused with a claim that commercial persistent agents universally have quadratic monetary cost. Real systems can use prompt-prefix caching, compaction, summarisation, context pruning, retrieval, server-side state, and changing provider architectures. Agent interfaces can also avoid replaying observations that are no longer useful. The theorem concerns cumulative information supplied under a specified naive replay model.

The bounded-$K$ assumption is equally important and potentially false. A more realistic model is:
\begin{equation}
K_i =
f(
|R_i|,
L_i,
D_i,
Q^{\mathrm{doc}}_i,
Q^{\mathrm{test}}_i,
Q^{\mathrm{arch}}_i,
Q^{\mathrm{retrieval}}_i
),
\label{eq:realistic-k}
\end{equation}
where $|R_i|$ is repository scale, $L_i$ task locality, $D_i$ dependency width, and the $Q$ terms represent documentation, test, architecture, and retrieval quality. The empirical question is not whether $K_i$ is literally constant. It is whether $K_i$ grows slowly enough, and with sufficiently high reconstruction fidelity, for externalised continuity to remain useful.

\subsection{Reconstruction cost as a falsification route}

The strongest obvious objection is that RaS merely moves cost from persistent context into repeated repository rediscovery. This is legitimate. RaS therefore defines reconstruction cost $C_{\mathrm{reconstruction}}$ as a first-class outcome rather than an implementation detail.

Future experiments should measure, where available:
\begin{itemize}
    \item input tokens attributable to reconstruction;
    \item files and bytes read;
    \item repository and symbol searches;
    \item history and dependency-navigation operations;
    \item tool calls and model calls;
    \item elapsed reconstruction time;
    \item cached versus uncached context;
    \item retries caused by incomplete or incorrect reconstruction.
\end{itemize}

We propose the \emph{Reconstruction Token Fraction}:
\begin{equation}
\mathrm{RTF}
=
\frac{T_{\mathrm{reconstruction}}}
{T_{\mathrm{RaS,input}}}.
\label{eq:rtf}
\end{equation}
RTF is meaningful only with a consistent classification rule. The boundary between “reconstruction” and “task reasoning” must be preregistered before outcomes are observed. Otherwise repository reading can be relabelled after the fact to make the architecture appear cheaper.

If reconstruction becomes excessively expensive as repositories grow, that weakens the RaS state-economics thesis even if forced-reset engineering success remains high.

\subsection{Provider-neutral cost decomposition}

For a persistent coding agent:
\begin{align}
C_{\mathrm{persistent}}
=&\ C_{\mathrm{reasoning}}
+ C_{\mathrm{context/history}}
+ C_{\mathrm{agent-state}} \nonumber\\
&+ C_{\mathrm{execution}}
+ C_{\mathrm{tooling}}
+ C_{\mathrm{orchestration}}.
\label{eq:cpersistent}
\end{align}

For RaS:
\begin{align}
C_{\mathrm{RaS}}
=&\ C_{\mathrm{reconstruction}}
+ C_{\mathrm{reasoning}}
+ C_{\mathrm{shallow-execution}} \nonumber\\
&+ C_{\mathrm{repository-state}}
+ C_{\mathrm{orchestration,RaS}}.
\label{eq:cras}
\end{align}

No component is assumed to be zero. The economic hypothesis is that for sufficiently long engineering programmes,
\begin{equation}
C_{\mathrm{RaS}} < C_{\mathrm{persistent}}
\end{equation}
while
\begin{equation}
Q_{\mathrm{engineering,RaS}}
\approx
Q_{\mathrm{engineering,persistent}}.
\end{equation}
The second condition is essential. Lower nominal model usage is not a saving if regressions, retries, or failed tasks increase sufficiently.

Three categories must remain separate:

\paragraph{Measured cost proxy.}
Observable model tokens, provider-reported usage, calls, files read, searches, elapsed time, execution operations, retries, and session persistence.

\paragraph{Theoretical infrastructure implication.}
A consequence derived from an explicit resource model, such as the possibility that lower persistent state per useful programme permits greater concurrency.

\paragraph{Unobservable provider-internal cost.}
GPU allocation, exact KV-cache cost, storage amortisation, scheduler behaviour, internal cache hit rates not exposed to the user, and provider margin.

The experiment cannot infer the third category from subscription pricing. Nor should a bundled chat product be described as “free inference”.

Table~\ref{tab:persistent-vs-ras} summarises the architectural comparison.

\input{tables/persistent-vs-ras}

Long-context serving research provides evidence that inference-state management can materially affect throughput. PagedAttention, for example, addresses KV-cache fragmentation and sharing in LLM serving \citep{kwon2023vllm}. That result motivates care around state burden but does not establish RaS savings. The defensible RaS claim remains conditional: if less high-capability historical state must persist per useful engineering programme, fixed inference infrastructure may potentially support more useful concurrent work.
